\documentclass[]{article}
\usepackage[T1]{fontenc}
\usepackage{lmodern}
\usepackage{amssymb,amsmath}
\usepackage{ifxetex,ifluatex}
\usepackage{fixltx2e} % provides \textsubscript
% use upquote if available, for straight quotes in verbatim environments
\IfFileExists{upquote.sty}{\usepackage{upquote}}{}
\ifnum 0\ifxetex 1\fi\ifluatex 1\fi=0 % if pdftex
  \usepackage[utf8]{inputenc}
\else % if luatex or xelatex
  \ifxetex
    \usepackage{mathspec}
    \usepackage{xltxtra,xunicode}
  \else
    \usepackage{fontspec}
  \fi
  \defaultfontfeatures{Mapping=tex-text,Scale=MatchLowercase}
  \newcommand{\euro}{€}
\fi
% use microtype if available
\IfFileExists{microtype.sty}{\usepackage{microtype}}{}
\ifxetex
  \usepackage[setpagesize=false, % page size defined by xetex
              unicode=false, % unicode breaks when used with xetex
              xetex]{hyperref}
\else
  \usepackage[unicode=true]{hyperref}
\fi
\hypersetup{breaklinks=true,
            bookmarks=true,
            pdfauthor={},
            pdftitle={},
            colorlinks=true,
            citecolor=blue,
            urlcolor=blue,
            linkcolor=magenta,
            pdfborder={0 0 0}}
\urlstyle{same}  % don't use monospace font for urls
\setlength{\parindent}{0pt}
\setlength{\parskip}{6pt plus 2pt minus 1pt}
\setlength{\emergencystretch}{3em}  % prevent overfull lines
\setcounter{secnumdepth}{0}

\author{}
\date{}
\usepackage{graphicx}
\usepackage{wrapfig}
\usepackage[figurename=Fig.,labelfont={it}]{caption}

\begin{document}

\begin{figure}[!h]
\centering
\includegraphics[width=0.5\textwidth]{test/test_fig1.pdf}\vspace{-10pt}%
\caption{\textbf{This is the caption for Figure 1.} It describes what Figure 1 shows, and perhaps how to interpret it. It may include details of the layout, or note landmarks and other things of interest.}
\label{Fig1}
\end{figure}

Fig. \ref{Fig1} is a good-looking figure, complete with a title, and
caption. The particular template used isn't much different than what you
get with standard pandoc figures, but it does show what is possible with
macros.

\begin{wrapfigure}{L}{0.3\textwidth}
\includegraphics[width=0.3\textwidth]{test/test_fig2.pdf}
\caption{\textbf{The Figure 2 caption.} This caption is different than that for Figure 1.}
\label{Fig2}
\end{wrapfigure}

Fig. \ref{Fig2} implements a wrapfig environment, and like Fig.
\ref{Fig1} it also has a title, and caption. The text of this paragraph
should be wrapping around the figure. These demonstration macros are
obviously latex, but html, or whatever else you need are just as easy to
implement. One reason to use macros is it allows you to neatly organize
everything in metadata or files as convenient, rather than spread out
across the document. It also allows you to re-use bits of text like
captions, which is very useful when needed. But mostly, it lets you
implement arbitrary formatting that sidesteps the limitations of pandoc.

You can also implement a simple file include macro. Here, text from the
file src/included\_tex.md is inserted:

file://src/included\_text.md

Notice that no ID is specified. It isn't needed unless you use
idvariables, which use the id to lookup the value to insert.

\end{document}
